% ============================================================
% APÉNDICE: FORMULACIONES MATEMÁTICAS Y ALGORITMOS
% ============================================================

\subsection{Formulaciones matemáticas y algoritmos de referencia}
\label{app:metodologia}

Este apéndice desarrolla las formulaciones y procedimientos algorítmicos asociados a las
metodologías discutidas en la Sección~\ref{sec:discusion_metodologica}. En primer lugar se
presenta una formulación MILP integrada que determina simultáneamente la secuencia semanal
de vuelos de las aeronaves, sus tiempos de operación y la asignación de la carga. Luego se
presentan formulaciones alternativas basadas en rotaciones, generación de columnas,
branch-and-price, descomposición de Benders y métodos matheurísticos.

La formulación principal utiliza únicamente el catálogo físico de tramos disponibles.
No se enumeran previamente vuelos con horarios específicos ni rotaciones completas. Para
cada aeronave se dispone de un número máximo de posiciones dentro de su secuencia semanal,
y el modelo decide qué tramo ocupa cada posición y en qué instante se realiza.


% ============================================================
\subsubsection{Red física y posiciones de vuelo}
\label{app:red_fisica}
% ============================================================

La red física se representa mediante un grafo dirigido

\[
G=(A,E),
\]

donde $A$ es el conjunto de aeropuertos y $E$ el conjunto de tramos dirigidos permitidos
entre ellos.

Para cada tramo $e\in E$ se conocen su aeropuerto de origen $o_e$, aeropuerto de destino
$d_e$ y tiempo de vuelo o block time $\tau_e$.

Para cada aeronave se define un conjunto de posiciones

\[
S=\{1,\ldots,\bar S\},
\]

donde $\bar S$ es una cota superior al número de vuelos que una aeronave puede efectuar
durante una semana. Esta cota puede obtenerse de manera conservadora a partir del horizonte
semanal y de la mínima duración posible de un ciclo vuelo--tiempo en tierra.

La utilización de $S$ no implica enumerar rutas. Por ejemplo, si

\[
y_{k,s,e}=1,
\]

el modelo está indicando que la aeronave $k$ realiza el tramo físico $e$ en la posición $s$
de su secuencia. La combinación completa de tramos de cada aeronave es determinada por el
solver.


% ============================================================
\subsubsection{Notación del modelo MILP}
\label{app:notacion_milp}
% ============================================================

\vspace{0.5em}
\noindent\textbf{Conjuntos.}

\begin{description}

    \item[$A$]
    conjunto de aeropuertos.

    \item[$E$]
    conjunto de tramos físicos dirigidos disponibles.

    \item[$K$]
    conjunto de aeronaves.

    \item[$S=\{1,\ldots,\bar S\}$]
    posiciones potenciales dentro de la secuencia semanal de una aeronave.

    \item[$S^-=\{1,\ldots,\bar S-1\}$]
    posiciones que pueden tener una operación posterior dentro de la misma semana.

    \item[$Q$]
    conjunto de demandas o commodities de carga.

    \item[$C$]
    conjunto de pares de tramos físicamente consecutivos,
    \[
    C=\{(e,e')\in E\times E:d_e=o_{e'}\}.
    \]

    \item[$L$]
    conjunto de mercados sujetos a frecuencia mínima.

    \item[$E_\ell\subseteq E$]
    conjunto de tramos que contribuyen al cumplimiento de la frecuencia del mercado
    $\ell\in L$.

    \item[$R_k^M$]
    conjunto de requerimientos de mantenimiento de la aeronave $k$.

\end{description}


\vspace{0.5em}
\noindent\textbf{Parámetros de red y aeronaves.}

\begin{description}

    \item[$O_{ie}$]
    parámetro igual a 1 si el tramo $e$ sale del aeropuerto $i$.

    \item[$D_{ie}$]
    parámetro igual a 1 si el tramo $e$ llega al aeropuerto $i$.

    \item[$\tau_e$]
    block time del tramo $e$.

    \item[$Q_{ke}$]
    capacidad efectiva de carga de la aeronave $k$ al operar el tramo $e$.

    \item[$c_{ke}^{\mathrm{op}}$]
    costo operacional de realizar el tramo $e$ con la aeronave $k$.

    \item[$R_{ke}\in\{0,1\}$]
    parámetro igual a 1 si el operador asociado a la aeronave $k$ posee los derechos de
    tráfico necesarios para operar el tramo $e$.

    \item[$A_{ke}\in\{0,1\}$]
    parámetro igual a 1 si el tramo $e$ es técnicamente factible para la aeronave $k$.

    \item[$T_{kee'}^{\mathrm{TAT}}$]
    tiempo mínimo en tierra requerido cuando la aeronave $k$ opera consecutivamente
    los tramos $e$ y $e'$.

    \item[$H$]
    duración del horizonte semanal, expresada en las mismas unidades de tiempo utilizadas
    en el modelo. Para una representación en horas, $H=168$.

    \item[$\Gamma$]
    constante suficientemente grande utilizada en restricciones Big-M.

\end{description}


\vspace{0.5em}
\noindent\textbf{Parámetros de demanda.}

Para cada $q\in Q$ se definen:

\begin{description}

    \item[$o_q$]
    aeropuerto de origen de la demanda.

    \item[$d_q$]
    aeropuerto de destino de la demanda.

    \item[$D_q$]
    demanda máxima disponible, medida en toneladas.

    \item[$r_q$]
    ingreso unitario por kilogramo transportado.

\end{description}

Se introduce además

\[
T^{\mathrm{tr}}\geq0,
\]

correspondiente al tiempo mínimo de conexión exigido cuando la carga cambia de aeronave.
Su valor numérico constituye un parámetro del modelo.

Para cada aeropuerto $i\in A$ se define

\[
\chi_i=
\begin{cases}
1, & \text{si está permitido realizar transbordos comerciales en }i,\\
0, & \text{en otro caso}.
\end{cases}
\]

Esto permite excluir, por ejemplo, eventos que correspondan exclusivamente a escalas
técnicas sin manipulación de carga.


\vspace{0.5em}
\noindent\textbf{Parámetros de frecuencia y mantenimiento.}

Para cada mercado $\ell\in L$ se define $F_\ell^{\min}$ como su frecuencia semanal mínima.

Para cada requerimiento $r\in R_k^M$ se conocen:

\[
a_{kr}^{M},
\qquad
\underline{t}_{kr}^{M},
\qquad
\overline{t}_{kr}^{M},
\qquad
d_{kr}^{M},
\]

correspondientes al aeropuerto de mantenimiento, inicio y fin de la ventana disponible, y
duración mínima de la actividad, respectivamente.


% ============================================================
\subsubsection{Variables de decisión}
\label{app:variables_milp}
% ============================================================

\vspace{0.5em}
\noindent\textbf{Variables asociadas a las aeronaves.}

\begin{align}
y_{kse}
&=
\begin{cases}
1, & \text{si la aeronave $k$ opera el tramo $e$ en la posición $s$},\\
0, & \text{en otro caso},
\end{cases}
\label{var:y}
\\
u_{ks}
&=
\begin{cases}
1, & \text{si la posición $s$ de la aeronave $k$ es utilizada},\\
0, & \text{en otro caso},
\end{cases}
\label{var:u}
\\
p_{ksee'}
&=
\begin{cases}
1, & \text{si $k$ realiza $e$ en $s$ y $e'$ en $s+1$},\\
0, & \text{en otro caso},
\end{cases}
\quad (e,e')\in C,
\label{var:p}
\\
\ell_{ks}
&=
\begin{cases}
1, & \text{si $s$ es la última posición utilizada por $k$},\\
0, & \text{en otro caso},
\end{cases}
\label{var:last}
\\
p^{W}_{ksee'}
&=
\begin{cases}
1, & \text{si $e$ es el último tramo de $k$ y $e'$ el primero}\\
   & \text{de la semana siguiente},
\end{cases}
\quad (e,e')\in C,
\label{var:wrap}
\\
t_{ks}^{\mathrm{dep}}
&\geq0,
\quad
\text{instante de salida de la posición $s$},
\label{var:tdep}
\\
t_{ks}^{\mathrm{arr}}
&\geq0,
\quad
\text{instante de llegada de la posición $s$}.
\label{var:tarr}
\end{align}


\vspace{0.5em}
\noindent\textbf{Variables asociadas a la carga.}

\begin{align}
x_{qks}
&\geq0,
\quad
\text{toneladas de $q$ transportadas en la posición $(k,s)$},
\label{var:x}
\\
b_{qks}
&\geq0,
\quad
\text{toneladas de $q$ que ingresan a la red en $(k,s)$},
\label{var:b}
\\
a_{qks}
&\geq0,
\quad
\text{toneladas de $q$ que terminan su recorrido en $(k,s)$},
\label{var:a}
\\
w_{qks}
&\geq0,
\quad
\text{toneladas de $q$ que continúan desde $(k,s)$ hacia $(k,s+1)$},
\label{var:w}
\\
z_{qksk's'}
&\geq0,
\quad
\text{toneladas de $q$ transferidas desde $(k,s)$ hacia $(k',s')$},
\quad k\neq k',
\label{var:z}
\\
g_{ksk's'}
&=
\begin{cases}
1, & \text{si existe una conexión de transbordo factible entre}\\
   & \text{$(k,s)$ y $(k',s')$},
\end{cases}
\quad k\neq k',
\label{var:g}
\\
s_q
&\geq0,
\quad
\text{toneladas totales de la demanda $q$ efectivamente servidas}.
\label{var:s}
\end{align}


\vspace{0.5em}
\noindent\textbf{Variables asociadas al mantenimiento.}

Se define

\[
m_{krs}=
\begin{cases}
1, & \text{si el mantenimiento $r$ de $k$ se realiza entre las posiciones $s$ y $s+1$},\\
0, & \text{en otro caso},
\end{cases}
\]

y

\[
\theta_{kr}\geq0,
\]

correspondiente al instante de inicio del mantenimiento $r$ de la aeronave $k$.


% ============================================================
\subsubsection{Formulación MILP integrada}
\label{app:modelo_milp}
% ============================================================

Para facilitar la lectura y permitir saltos de página, la formulación se presenta en tres
bloques: programación de aeronaves, flujo de carga y restricciones operacionales. Los tres
bloques constituyen conjuntamente un único MILP.


% ------------------------------------------------------------
% FUNCIÓN OBJETIVO
% ------------------------------------------------------------

\vspace{0.5em}
\noindent\textbf{Función objetivo.}

La función objetivo maximiza el margen neto semanal:

\begin{align}
\max \quad Z
={}&
\sum_{q\in Q}1000\,r_qs_q
-
\sum_{k\in K}\sum_{s\in S}\sum_{e\in E}
c_{ke}^{\mathrm{op}}y_{kse}
-
C^{\mathrm{cargo}}(b,a,z).
\tag{FO}
\label{eq:fo}
\end{align}

La función $C^{\mathrm{cargo}}(b,a,z)$ representa los costos lineales de manipulación,
embarque, descarga y transbordo de la carga de acuerdo con las tarifas correspondientes a
cada aeropuerto. Esta notación permite mantener legible la formulación sin alterar su
estructura lineal.


% ------------------------------------------------------------
% BLOQUE A
% ------------------------------------------------------------

\vspace{0.5em}
\noindent\textbf{Bloque A: selección, secuencia y tiempo de las aeronaves.}

Cada posición utilizada contiene exactamente un tramo:

\begin{align}
\sum_{e\in E}y_{kse}
&=
u_{ks},
\quad
\forall k\in K,\ s\in S.
\tag{R1}
\label{eq:r1}
\end{align}

Las posiciones se utilizan consecutivamente desde el comienzo de la secuencia:

\begin{align}
u_{k,s+1}
&\leq
u_{ks},
\quad
\forall k\in K,\ s\in S^-.
\tag{R2a}
\label{eq:r2a}
\\
\ell_{ks}
&=
u_{ks}-u_{k,s+1},
\quad
\forall k\in K,\ s\in S^-,
\tag{R2b}
\label{eq:r2b}
\\
\ell_{k,\bar S}
&=
u_{k,\bar S},
\quad
\forall k\in K.
\tag{R2c}
\label{eq:r2c}
\end{align}

Para dos posiciones consecutivas se selecciona un par de tramos físicamente compatible:

\begin{align}
p_{ksee'}
&\leq y_{kse},
\quad
\forall k,\ s\in S^-,\ (e,e')\in C,
\tag{R3a}
\label{eq:r3a}
\\
p_{ksee'}
&\leq y_{k,s+1,e'},
\quad
\forall k,\ s\in S^-,\ (e,e')\in C,
\tag{R3b}
\label{eq:r3b}
\\
p_{ksee'}
&\geq
y_{kse}+y_{k,s+1,e'}-1,
\quad
\forall k,\ s\in S^-,\ (e,e')\in C,
\tag{R3c}
\label{eq:r3c}
\\
\sum_{(e,e')\in C}p_{ksee'}
&=
u_{k,s+1},
\quad
\forall k\in K,\ s\in S^-.
\tag{R3d}
\label{eq:r3d}
\end{align}

Los tiempos de llegada se obtienen a partir de los tiempos de salida y del tramo seleccionado:

\begin{align}
t_{ks}^{\mathrm{arr}}
&=
t_{ks}^{\mathrm{dep}}
+
\sum_{e\in E}\tau_e y_{kse},
\quad
\forall k\in K,\ s\in S.
\tag{R4a}
\label{eq:r4a}
\end{align}

Entre vuelos consecutivos debe cumplirse el tiempo mínimo en tierra:

\begin{align}
t_{k,s+1}^{\mathrm{dep}}
&\geq
t_{ks}^{\mathrm{arr}}
+
\sum_{(e,e')\in C}
T_{kee'}^{\mathrm{TAT}}p_{ksee'}
-
\Gamma(1-u_{k,s+1}),
\nonumber\\[-1mm]
&\hspace{5cm}
\forall k\in K,\ s\in S^-.
\tag{R4b}
\label{eq:r4b}
\end{align}

Los tiempos quedan dentro del horizonte semanal cuando la posición es utilizada:

\begin{align}
0
\leq t_{ks}^{\mathrm{dep}}
&\leq
H\,u_{ks},
\quad
\forall k\in K,\ s\in S,
\tag{R4c}
\label{eq:r4c}
\\
0
\leq t_{ks}^{\mathrm{arr}}
&\leq
H\,u_{ks},
\quad
\forall k\in K,\ s\in S.
\tag{R4d}
\label{eq:r4d}
\end{align}

El último tramo debe conectarse físicamente con el primero para formar una rotación semanal
cíclica:

\begin{align}
p^W_{ksee'}
&\leq
\ell_{ks},
\quad
\forall k,\ s,\ (e,e')\in C,
\tag{R5a}
\label{eq:r5a}
\\
p^W_{ksee'}
&\leq
y_{kse},
\quad
\forall k,\ s,\ (e,e')\in C,
\tag{R5b}
\label{eq:r5b}
\\
p^W_{ksee'}
&\leq
y_{k,1,e'},
\quad
\forall k,\ s,\ (e,e')\in C,
\tag{R5c}
\label{eq:r5c}
\\
p^W_{ksee'}
&\geq
\ell_{ks}+y_{kse}+y_{k,1,e'}-2,
\quad
\forall k,\ s,\ (e,e')\in C,
\tag{R5d}
\label{eq:r5d}
\\
\sum_{s\in S}\sum_{(e,e')\in C}
p^W_{ksee'}
&=
u_{k1},
\quad
\forall k\in K.
\tag{R5e}
\label{eq:r5e}
\end{align}

La conexión entre el último vuelo y el primero de la semana siguiente debe ser también
temporalmente factible:

\begin{align}
t_{k1}^{\mathrm{dep}}+H
&\geq
t_{ks}^{\mathrm{arr}}
+
T_{kee'}^{\mathrm{TAT}}
-
\Gamma(1-p^W_{ksee'}),
\nonumber\\[-1mm]
&\hspace{3.8cm}
\forall k,\ s,\ (e,e')\in C.
\tag{R5f}
\label{eq:r5f}
\end{align}


% ------------------------------------------------------------
% BLOQUE B
% ------------------------------------------------------------

\vspace{0.5em}
\noindent\textbf{Bloque B: flujo multicommodity de carga.}

La carga transportada en una operación está limitada por la capacidad de la aeronave y el
tramo seleccionado:

\begin{align}
\sum_{q\in Q}x_{qks}
&\leq
\sum_{e\in E}Q_{ke}y_{kse},
\quad
\forall k\in K,\ s\in S.
\tag{R6}
\label{eq:r6}
\end{align}

La conservación del flujo al entrar a una operación se expresa mediante

\begin{align}
x_{qks}
={}&
b_{qks}
+
w_{q,k,s-1}
+
\sum_{\substack{k'\in K\\k'\neq k}}
\sum_{s'\in S}
z_{qk's'ks},
\nonumber\\[-1mm]
&\hspace{4cm}
\forall q\in Q,\ k\in K,\ s\in S,
\tag{R7a}
\label{eq:r7a}
\end{align}

donde se adopta la convención

\[
w_{q,k,0}=0.
\]

La conservación al abandonar una operación es

\begin{align}
x_{qks}
={}&
a_{qks}
+
w_{qks}
+
\sum_{\substack{k'\in K\\k'\neq k}}
\sum_{s'\in S}
z_{qksk's'},
\nonumber\\[-1mm]
&\hspace{4cm}
\forall q\in Q,\ k\in K,\ s\in S,
\tag{R7b}
\label{eq:r7b}
\end{align}

con

\[
w_{q,k,\bar S}=0.
\]

La carga puede continuar en la misma aeronave entre posiciones consecutivas:

\begin{align}
w_{qks}
&\leq
D_q\,u_{k,s+1},
\quad
\forall q\in Q,\ k\in K,\ s\in S^-.
\tag{R8a}
\label{eq:r8a}
\end{align}

Cuando existe cambio de aeronave se utiliza la variable de conexión $g_{ksk's'}$. Esta solo
puede activarse si ambas operaciones existen:

\begin{align}
g_{ksk's'}
&\leq
u_{ks},
\quad
\forall k\neq k',\ s,s'\in S,
\tag{R8b}
\label{eq:r8b}
\\
g_{ksk's'}
&\leq
u_{k's'},
\quad
\forall k\neq k',\ s,s'\in S.
\tag{R8c}
\label{eq:r8c}
\end{align}

Para que la conexión exista, el destino de la primera operación debe coincidir con el origen
de la segunda. Para cada aeropuerto $i\in A$:

\begin{align}
\sum_{e\in E}D_{ie}y_{kse}
-
\sum_{e\in E}O_{ie}y_{k's'e}
&\leq
1-g_{ksk's'},
\nonumber\\[-1mm]
&\hspace{2.7cm}
\forall i,\ k\neq k',\ s,s',
\tag{R8d}
\label{eq:r8d}
\\
\sum_{e\in E}O_{ie}y_{k's'e}
-
\sum_{e\in E}D_{ie}y_{kse}
&\leq
1-g_{ksk's'},
\nonumber\\[-1mm]
&\hspace{2.7cm}
\forall i,\ k\neq k',\ s,s'.
\tag{R8e}
\label{eq:r8e}
\end{align}

La conexión solo puede utilizar un aeropuerto donde se permita manipular carga:

\begin{align}
g_{ksk's'}
&\leq
1-
\sum_{e\in E}D_{ie}y_{kse}
+
\chi_i,
\nonumber\\[-1mm]
&\hspace{2.7cm}
\forall i,\ k\neq k',\ s,s'.
\tag{R8f}
\label{eq:r8f}
\end{align}

Cuando se cambia de aeronave debe respetarse el tiempo mínimo de transbordo:

\begin{align}
t_{k's'}^{\mathrm{dep}}
&\geq
t_{ks}^{\mathrm{arr}}
+
T^{\mathrm{tr}}
-
\Gamma(1-g_{ksk's'}),
\nonumber\\[-1mm]
&\hspace{3.2cm}
\forall k\neq k',\ s,s'.
\tag{R8g}
\label{eq:r8g}
\end{align}

El flujo transferido solo puede ser positivo cuando dicha conexión está habilitada:

\begin{align}
z_{qksk's'}
&\leq
D_q\,g_{ksk's'},
\quad
\forall q\in Q,\ k\neq k',\ s,s'.
\tag{R8h}
\label{eq:r8h}
\end{align}

La carga solo puede comenzar su recorrido en su aeropuerto de origen:

\begin{align}
b_{qks}
&\leq
D_q
\sum_{\substack{e\in E\\o_e=o_q}}
y_{kse},
\quad
\forall q\in Q,\ k\in K,\ s\in S.
\tag{R9a}
\label{eq:r9a}
\end{align}

Análogamente, solo puede finalizarlo en su aeropuerto de destino:

\begin{align}
a_{qks}
&\leq
D_q
\sum_{\substack{e\in E\\d_e=d_q}}
y_{kse},
\quad
\forall q\in Q,\ k\in K,\ s\in S.
\tag{R9b}
\label{eq:r9b}
\end{align}

La cantidad total servida debe coincidir tanto con la carga que ingresa como con la que
alcanza su destino:

\begin{align}
s_q
&=
\sum_{k\in K}\sum_{s\in S}b_{qks},
\quad
\forall q\in Q,
\tag{R9c}
\label{eq:r9c}
\\
s_q
&=
\sum_{k\in K}\sum_{s\in S}a_{qks},
\quad
\forall q\in Q,
\tag{R9d}
\label{eq:r9d}
\\
0\leq s_q
&\leq
D_q,
\quad
\forall q\in Q.
\tag{R9e}
\label{eq:r9e}
\end{align}


% ------------------------------------------------------------
% BLOQUE C
% ------------------------------------------------------------

\vspace{0.5em}
\noindent\textbf{Bloque C: restricciones operacionales adicionales.}

Los compromisos comerciales de frecuencia se expresan como

\begin{align}
\sum_{k\in K}\sum_{s\in S}
\sum_{e\in E_\ell}
y_{kse}
&\geq
F_\ell^{\min},
\quad
\forall \ell\in L.
\tag{R10}
\label{eq:r10}
\end{align}

Los derechos de tráfico del operador asociado a cada aeronave se imponen mediante

\begin{align}
y_{kse}
&\leq
R_{ke},
\quad
\forall k\in K,\ s\in S,\ e\in E.
\tag{R11}
\label{eq:r11}
\end{align}

La factibilidad técnica de cada tramo se impone mediante

\begin{align}
y_{kse}
&\leq
A_{ke},
\quad
\forall k\in K,\ s\in S,\ e\in E.
\tag{R12}
\label{eq:r12}
\end{align}

Cada requerimiento de mantenimiento debe ser asignado exactamente una vez a un intervalo
entre vuelos consecutivos:

\begin{align}
\sum_{s\in S^-}m_{krs}
&=
1,
\quad
\forall k\in K,\ r\in R_k^M.
\tag{R13a}
\label{eq:r13a}
\end{align}

El mantenimiento solo puede asignarse a un intervalo que realmente exista:

\begin{align}
m_{krs}
&\leq
u_{k,s+1},
\quad
\forall k,\ r\in R_k^M,\ s\in S^-.
\tag{R13b}
\label{eq:r13b}
\end{align}

La aeronave debe encontrarse en el aeropuerto de mantenimiento al terminar el vuelo anterior:

\begin{align}
m_{krs}
&\leq
\sum_{\substack{e\in E\\d_e=a_{kr}^{M}}}
y_{kse},
\quad
\forall k,\ r\in R_k^M,\ s\in S^-.
\tag{R13c}
\label{eq:r13c}
\end{align}

Asimismo, el siguiente vuelo debe partir desde dicho aeropuerto:

\begin{align}
m_{krs}
&\leq
\sum_{\substack{e\in E\\o_e=a_{kr}^{M}}}
y_{k,s+1,e},
\quad
\forall k,\ r\in R_k^M,\ s\in S^-.
\tag{R13d}
\label{eq:r13d}
\end{align}

El inicio del mantenimiento debe ubicarse dentro de su ventana temporal:

\begin{align}
\underline{t}_{kr}^{M}
\leq
\theta_{kr}
&\leq
\overline{t}_{kr}^{M}-d_{kr}^{M},
\quad
\forall k,\ r\in R_k^M.
\tag{R13e}
\label{eq:r13e}
\end{align}

Cuando $m_{krs}=1$, el mantenimiento debe comenzar después de la llegada y terminar antes del
siguiente vuelo:

\begin{align}
\theta_{kr}
&\geq
t_{ks}^{\mathrm{arr}}
-
\Gamma(1-m_{krs}),
\nonumber\\[-1mm]
&\hspace{3cm}
\forall k,\ r\in R_k^M,\ s\in S^-,
\tag{R13f}
\label{eq:r13f}
\\
\theta_{kr}+d_{kr}^{M}
&\leq
t_{k,s+1}^{\mathrm{dep}}
+
\Gamma(1-m_{krs}),
\nonumber\\[-1mm]
&\hspace{3cm}
\forall k,\ r\in R_k^M,\ s\in S^-.
\tag{R13g}
\label{eq:r13g}
\end{align}

Finalmente, los dominios de las variables son

\begin{align}
y_{kse},u_{ks},p_{ksee'},\ell_{ks},p^W_{ksee'},
m_{krs},g_{ksk's'}
&\in
\{0,1\},
\tag{R14a}
\label{eq:r14a}
\\
t_{ks}^{\mathrm{dep}},t_{ks}^{\mathrm{arr}},
\theta_{kr},
x_{qks},b_{qks},a_{qks},w_{qks},
z_{qksk's'},s_q
&\geq
0.
\tag{R14b}
\label{eq:r14b}
\end{align}


% ============================================================
\subsubsection{Interpretación del modelo MILP}
\label{app:interpretacion_milp}
% ============================================================

\vspace{0.5em}
\noindent\textbf{Variables $y_{kse}$ y $u_{ks}$.}

La variable $y_{kse}$ determina conjuntamente qué aeronave opera cada tramo y en qué posición
de su secuencia semanal lo realiza. No existe un conjunto de vuelos previamente calendarizados:
la hora de operación se determina mediante $t_{ks}^{\mathrm{dep}}$ y
$t_{ks}^{\mathrm{arr}}$.

La variable $u_{ks}$ indica si la posición $s$ está ocupada. R1 obliga a seleccionar
exactamente un tramo cuando la posición se utiliza.


\vspace{0.5em}
\noindent\textbf{R2. Estructura de las posiciones.}

Las restricciones R2a--R2c evitan posiciones vacías en medio de una rotación. Si una posición
$s+1$ es utilizada, entonces $s$ también debe serlo.

La variable $\ell_{ks}$ identifica automáticamente cuál es la última posición utilizada, lo
que permite imponer posteriormente el cierre semanal sin conocer de antemano cuántos vuelos
realizará cada aeronave.


\vspace{0.5em}
\noindent\textbf{R3. Continuidad espacial.}

Las variables $p_{ksee'}$ vinculan dos posiciones consecutivas.

Como solo se definen para

\[
(e,e')\in C,
\]

y por definición

\[
d_e=o_{e'},
\]

R3 garantiza que el aeropuerto donde termina una operación coincida con aquel donde comienza
la siguiente.

Por lo tanto, la ruta completa de una aeronave surge de las decisiones del modelo:

\[
e_1\rightarrow e_2\rightarrow e_3\rightarrow\cdots,
\]

sin enumerar previamente ninguna rotación.


\vspace{0.5em}
\noindent\textbf{R4. Dimensión temporal.}

R4a calcula la llegada de cada operación a partir del horario de salida y del block time del
tramo seleccionado.

R4b impone el tiempo mínimo en tierra entre vuelos consecutivos. Como las horas de salida son
variables continuas, el modelo puede decidir libremente cuándo realizar cada vuelo dentro del
horizonte semanal.

La formulación evita así discretizar el tiempo generando, por ejemplo, copias de un mismo
tramo a las 8:00, 9:00, 10:00, etc.


\vspace{0.5em}
\noindent\textbf{R5. Periodicidad semanal.}

Las restricciones R5a--R5e conectan la última posición utilizada con la primera mediante un
par de tramos físicamente compatible.

R5f impone además que exista suficiente tiempo entre la última llegada y la primera salida
de la semana siguiente:

\[
t_{k1}^{\mathrm{dep}}+H
\geq
t_{ks}^{\mathrm{arr}}
+
T_{kee'}^{\mathrm{TAT}}.
\]

La secuencia de cada aeronave forma así una rotación semanal repetible.


\vspace{0.5em}
\noindent\textbf{R6. Capacidad y acoplamiento aeronave--carga.}

R6 constituye el principal vínculo entre ambos componentes:

\[
\sum_{q\in Q}x_{qks}
\leq
\sum_{e\in E}Q_{ke}y_{kse}.
\]

La capacidad de la posición $(k,s)$ depende del tramo efectivamente seleccionado.

Si no existe una operación en dicha posición,

\[
u_{ks}=0,
\]

entonces R1 implica que todos los $y_{kse}$ son cero y, por R6, ninguna carga puede utilizar
esa operación.


\vspace{0.5em}
\noindent\textbf{R7. Conservación del flujo de carga.}

R7a explica de dónde proviene la carga presente en una determinada operación. Puede haber
sido embarcada inicialmente, haber llegado mediante un transbordo desde otra aeronave o haber
continuado en el mismo avión desde la posición anterior.

R7b describe los tres destinos posibles después del vuelo: descarga final, transbordo hacia
otra aeronave o continuación en el mismo avión.

De este modo, el itinerario completo de cada commodity también surge de las decisiones del
modelo y no necesita ser enumerado previamente.


\vspace{0.5em}
\noindent\textbf{R8. Conexiones y transbordos.}

R8a representa una conexión \textit{through}: la carga permanece en el mismo avión y continúa
hacia su siguiente operación.

Las restricciones R8b--R8f determinan si dos operaciones de aeronaves distintas pueden
conectarse físicamente en un mismo aeropuerto.

R8g incorpora la principal condición temporal de transbordo:

\[
t_{k's'}^{\mathrm{dep}}
-
t_{ks}^{\mathrm{arr}}
\geq
T^{\mathrm{tr}}.
\]

Por lo tanto, una carga puede cambiar de avión en cualquier aeropuerto donde se permita
manipular carga, siempre que exista suficiente tiempo entre ambos vuelos.

R8h impide enviar carga por una conexión que no haya sido habilitada.


\vspace{0.5em}
\noindent\textbf{R9. Origen, destino y demanda disponible.}

R9a permite que una demanda entre a la red únicamente mediante una operación que parte desde
su aeropuerto de origen.

R9b exige que solo pueda finalizar su recorrido en una operación que llegue a su destino.

R9c y R9d igualan la cantidad que entra a la red con la cantidad efectivamente entregada,
mientras R9e limita el servicio a la demanda máxima disponible:

\[
s_q\leq D_q.
\]

No existe obligación de atender toda la demanda.


\vspace{0.5em}
\noindent\textbf{R10. Frecuencias mínimas.}

R10 exige que los mercados sujetos a compromisos de frecuencia reciban al menos el número
mínimo de operaciones establecido.


\vspace{0.5em}
\noindent\textbf{R11. Derechos de tráfico.}

R11 impide seleccionar combinaciones aeronave--tramo que no estén autorizadas para el operador
correspondiente.


\vspace{0.5em}
\noindent\textbf{R12. Factibilidad técnica.}

R12 elimina tramos incompatibles con una aeronave considerando las restricciones técnicas
aplicables.

El parámetro $A_{ke}$ puede incorporar, entre otras condiciones, duración máxima de vuelo
continuo, compatibilidad aeroportuaria y necesidad de realizar escalas técnicas.

Cuando un viaje requiere una escala técnica, el modelo puede construirla seleccionando dos o
más tramos consecutivos del catálogo físico. No es necesario representar el viaje completo
como una alternativa enumerada.


\vspace{0.5em}
\noindent\textbf{R13. Mantenimiento.}

R13 asigna cada mantenimiento a uno de los períodos de tierra existentes entre vuelos
consecutivos.

Las restricciones de ubicación aseguran que la aeronave permanezca en el aeropuerto correcto,
mientras que R13e--R13g verifican que la actividad pueda completarse dentro de su ventana
temporal.


\vspace{0.5em}
\noindent\textbf{R14. Dominio de las variables.}

R14 establece la naturaleza binaria de las decisiones de programación y conexión y la no
negatividad de las variables continuas de tiempo y carga.


% ============================================================
\subsubsection{Tratamiento de reglas operacionales adicionales}
\label{app:reglas_adicionales}
% ============================================================

\vspace{0.5em}
\noindent\textbf{Escalas técnicas sin manipulación de carga.}

Una escala técnica puede ser parte de la secuencia de una aeronave mediante dos tramos
consecutivos. Si el aeropuerto solo permite abastecimiento de combustible, se utiliza

\[
\chi_i=0,
\]

impidiendo transbordos entre aeronaves en dicho punto.

La carga puede permanecer en el mismo avión mediante $w_{qks}$, lo que permite representar
una escala técnica sin descarga ni transferencia.


\vspace{0.5em}
\noindent\textbf{Interacción combustible--payload.}

La capacidad $Q_{ke}$ puede depender de la longitud del tramo y de las condiciones de
combustible.

Mientras no se disponga de una curva completa combustible--payload, el modelo utiliza la
capacidad efectiva que pueda determinarse a partir de la información disponible, sin imponer
una relación funcional no especificada.


\vspace{0.5em}
\noindent\textbf{Tonelaje mínimo asociado a escalas adicionales.}

La regla operacional que exige un tonelaje mínimo cuando se incorpora una escala adicional
requiere identificar con precisión qué secuencias son consideradas como tales y qué cantidad
de carga debe utilizarse para verificar el umbral.

Una vez definida dicha clasificación, puede introducirse un conjunto $C^{L}\subseteq C$ de
transiciones sujetas a la regla y una variable binaria que distinga entre una operación ferry
y una operación comercial. La estructura genérica es

\[
L_{ee'}^{\min}\,\nu_{ksee'}
\leq
X_{ksee'}
\leq
\overline{Q}_{k}\,\nu_{ksee'},
\]

donde $X_{ksee'}$ representa el tonelaje relevante para la regla y
$\nu_{ksee'}=0$ permite la operación ferry sin carga.

La definición exacta de $X_{ksee'}$ debe seguir la interpretación operacional entregada para
esta condición, evitando imponer una definición adicional no contenida en los datos.


\vspace{0.5em}
\noindent\textbf{Demanda diaria.}

Cuando la demanda se desagrega por día, cada commodity incorpora también su instante o ventana
de disponibilidad.

Las variables $b_{qks}$ pueden vincularse a esta ventana mediante restricciones Big-M sobre
$t_{ks}^{\mathrm{dep}}$. De esta forma, una carga solo puede comenzar su recorrido cuando se
encuentra disponible y no puede trasladarse artificialmente entre días.


% ============================================================
\subsubsection{Estructura general del MILP}
\label{app:estructura_milp}
% ============================================================

La formulación puede interpretarse como dos sistemas acoplados.

El primero construye cada rotación:

\[
\boxed{
\begin{array}{c}
\text{Tramos físicos}\\
E
\end{array}
}
\quad\longrightarrow\quad
\boxed{
\begin{array}{c}
y_{kse}\\
\text{selección de tramo}
\end{array}
}
\quad\longrightarrow\quad
\boxed{
\begin{array}{c}
t_{ks}^{\mathrm{dep}},\,t_{ks}^{\mathrm{arr}}\\
\text{programación temporal}
\end{array}
}
\]

El segundo distribuye los commodities sobre las operaciones resultantes:

\[
\boxed{
\begin{array}{c}
x_{qks}\\
\text{carga por vuelo}
\end{array}
}
\quad\longrightarrow\quad
\boxed{
\begin{array}{c}
w_{qks}\\
\text{conexión through}
\end{array}
}
\quad\text{o}\quad
\boxed{
\begin{array}{c}
z_{qksk's'}\\
\text{transbordo}
\end{array}
}
\]

Ambas capas se vinculan principalmente mediante R6:

\[
\boxed{
\sum_{q\in Q}x_{qks}
\leq
\sum_{e\in E}Q_{ke}y_{kse}
}
\]

Esta restricción expresa directamente que la carga solo puede utilizar capacidad generada por
una aeronave que efectivamente realiza un tramo.


% ============================================================
\subsubsection{Formulación alternativa basada en rotaciones}
\label{app:route_formulation}
% ============================================================

La formulación MILP utiliza decisiones por posición y tramo. Una representación alternativa
consiste en definir variables asociadas a rotaciones semanales completas.

Sea $R_k$ el conjunto conceptual de todas las rotaciones factibles de la aeronave $k$ y

\begin{equation}
\lambda_{kr}
=
\begin{cases}
1, & \text{si la aeronave $k$ utiliza la rotación $r$},\\
0, & \text{en otro caso}.
\end{cases}
\label{eq:lambda}
\end{equation}

Cada aeronave puede seleccionar como máximo una rotación:

\begin{equation}
\sum_{r\in R_k}\lambda_{kr}
\leq1,
\quad
\forall k\in K.
\label{eq:one_rotation}
\end{equation}

Una rotación contiene de antemano una secuencia espacial y temporalmente factible de vuelos,
incluyendo las restricciones propias de la aeronave.

El problema es que el número de elementos de $R_k$ puede crecer exponencialmente. Por ello,
esta formulación no requiere enumerar previamente $R_k$ cuando se utiliza generación de
columnas.


% ============================================================
\subsubsection{Generación de columnas}
\label{app:column_generation}
% ============================================================

La generación de columnas comienza con un subconjunto reducido

\[
\overline{R}_k\subset R_k.
\]

Estas rotaciones forman un Restricted Master Problem (RMP). Una vez resuelta su relajación
lineal, los precios duales permiten determinar si existe alguna rotación todavía no incluida
que pueda mejorar la solución.

En una formulación de minimización, el costo reducido genérico puede representarse mediante

\begin{equation}
\bar c_{kr}
=
c_{kr}
-
\alpha_k
-
\sum_m\pi_m a_{mkr},
\label{eq:reduced_cost}
\end{equation}

donde $\alpha_k$ y $\pi_m$ son variables duales del RMP.

El problema de \textit{pricing} busca

\begin{equation}
\min_{r\in R_k}\bar c_{kr}.
\label{eq:pricing}
\end{equation}

Si encuentra una rotación con costo reducido favorable, esta se incorpora al RMP y el
procedimiento se repite.

Esta estructura es consistente con los enfoques desarrollados por
\shortciteA{Derigs2009}, \shortciteA{Liang2015},
\citeA{DerigsFriederichs2013}, \shortciteA{Xiao2022} y
\shortciteA{Zhu2026}.


\vspace{0.5em}
\noindent\textbf{Algoritmo conceptual de generación de columnas.}

\begin{enumerate}

    \item Construir un conjunto inicial reducido de rotaciones factibles.

    \item Resolver la relajación lineal del Restricted Master Problem.

    \item Obtener los valores duales de las restricciones del RMP.

    \item Resolver el problema de \textit{pricing} para cada aeronave o tipo de aeronave.

    \item Incorporar las rotaciones con costo reducido favorable.

    \item Resolver nuevamente el RMP.

    \item Repetir hasta que ningún problema de \textit{pricing} encuentre una columna
    mejorante.

    \item Construir posteriormente una solución entera utilizando las columnas obtenidas.

\end{enumerate}


% ============================================================
\subsubsection{Branch-and-price y branch-and-price-and-cut}
\label{app:branch_price}
% ============================================================

La generación de columnas obtiene el óptimo de la relajación lineal del problema maestro,
pero las variables $\lambda_{kr}$ pueden resultar fraccionales.

Branch-and-price incorpora generación de columnas dentro de un árbol de branch-and-bound:

\[
\boxed{
\text{Branch-and-Bound}
+
\text{Column Generation}
=
\text{Branch-and-Price}.
}
\]

Cuando además se incorporan dinámicamente desigualdades válidas,

\[
\boxed{
\text{Branch-and-Price}
+
\text{Cuts}
=
\text{Branch-and-Price-and-Cut}.
}
\]

La metodología general de branch-and-price es desarrollada por
\shortciteA{Barnhart1998}, mientras que \citeA{DerigsFriederichs2013} aplican
branch-and-price-and-cut a planificación integrada de carga aérea.


% ============================================================
\subsubsection{Descomposición de Benders}
\label{app:benders}
% ============================================================

El MILP presenta una separación natural entre las decisiones discretas de programación de
aeronaves y las decisiones continuas de flujo de carga.

Sea $m$ el vector que agrupa las decisiones de programación:

\[
m=(y,u,p,\ell,p^W,t,m^M),
\]

donde $m^M$ representa las decisiones de mantenimiento.

Fijada una programación, quedan determinados los vuelos, sus horarios, sus capacidades y las
conexiones disponibles. El problema restante consiste en distribuir la carga mediante las
variables continuas

\[
h=(x,b,a,w,z,s).
\]

La estructura puede representarse abstractamente como

\begin{equation}
\max_{m\in Y}
\left\{
-c^\top m+Q(m)
\right\},
\label{eq:benders_outer}
\end{equation}

donde $Q(m)$ representa el beneficio máximo asociado al flujo de carga para una programación
determinada.

El subproblema puede expresarse genéricamente como

\begin{align}
Q(m)=\max_h\quad
& r^\top h
\label{eq:benders_sub_obj}
\\
\text{s.a.}\quad
& Bh\leq d,
\label{eq:benders_sub1}
\\
& Dh\leq Gm,
\label{eq:benders_sub2}
\\
& h\geq0.
\label{eq:benders_sub3}
\end{align}

Las restricciones \eqref{eq:benders_sub1} representan las condiciones propias del flujo de
carga, mientras que \eqref{eq:benders_sub2} vinculan dicho flujo con la capacidad y conexiones
generadas por el problema maestro.

El dual puede representarse mediante

\begin{align}
Q(m)=\min_{\pi,\mu}\quad
& d^\top\pi+(Gm)^\top\mu
\label{eq:benders_dual_obj}
\\
\text{s.a.}\quad
& B^\top\pi+D^\top\mu\geq r,
\label{eq:benders_dual_cons}
\\
& \pi,\mu\geq0.
\end{align}

Introduciendo una variable $\theta$ que aproxima el beneficio del subproblema, cada solución
dual extrema permite construir un corte

\begin{equation}
\theta
\leq
d^\top\pi^r
+
(Gm)^\top\mu^r.
\label{eq:benders_cut}
\end{equation}

Esta estructura es cercana al enfoque utilizado por \shortciteA{Li2006} para integrar
\textit{fleet assignment} y \textit{cargo routing}.


\vspace{0.5em}
\noindent\textbf{Algoritmo conceptual de Benders.}

\begin{enumerate}

    \item Resolver el problema maestro y obtener una programación de aeronaves.

    \item Fijar los vuelos, horarios, capacidades y conexiones generados por dicha
    programación.

    \item Resolver el subproblema continuo de flujo de carga.

    \item Obtener las variables duales del subproblema.

    \item Generar un corte de Benders y agregarlo al problema maestro.

    \item Actualizar las cotas del problema.

    \item Repetir hasta alcanzar la tolerancia de optimalidad o el límite de tiempo.

\end{enumerate}


% ============================================================
\subsubsection{MIP restringido y matheurísticas}
\label{app:restricted_mip}
% ============================================================

Una alternativa práctica a branch-and-price consiste en utilizar generación de columnas para
obtener un conjunto reducido

\[
\overline{R}_k\subset R_k
\]

y posteriormente resolver un MILP imponiendo integralidad sobre dichas columnas:

\begin{equation}
\max
\left\{
Z(\lambda,h):
\lambda_{kr}\in\{0,1\},
\quad
r\in\overline{R}_k
\right\}.
\label{eq:restricted_mip}
\end{equation}

Si todas las columnas corresponden a rotaciones factibles, la solución obtenida también es
factible para el problema original. No existe, sin embargo, garantía de optimalidad global
porque pueden existir rotaciones relevantes que no hayan sido generadas.

Estrategias de \textit{diving}, fix-and-optimize y reoptimización pueden utilizarse para
mejorar el incumbente. Este tipo de metodología es consistente con los procedimientos
desarrollados por \shortciteA{Xiao2022} y \shortciteA{Zhu2026}.


% ============================================================
\subsubsection{Adaptive Large Neighborhood Search}
\label{app:alns}
% ============================================================

Adaptive Large Neighborhood Search (ALNS) modifica iterativamente partes significativas de
una solución.

Si $S^{\mathrm{sol}}$ representa la solución actual, una iteración puede escribirse como

\[
S^{\mathrm{sol}}
\longrightarrow
S^{-}
\longrightarrow
S',
\]

donde $S^{-}$ corresponde a una solución parcialmente destruida y $S'$ a una solución
reconstruida.

El procedimiento general consiste en:

\begin{enumerate}

    \item Construir una solución inicial.

    \item Seleccionar adaptativamente un operador de destrucción.

    \item Liberar una parte de las decisiones de programación o flujo.

    \item Reconstruirlas mediante una heurística o un problema de optimización restringido.

    \item Aceptar o rechazar la nueva solución de acuerdo con el criterio utilizado.

    \item Actualizar la valoración de los operadores.

    \item Repetir hasta alcanzar el criterio de término.

\end{enumerate}

\shortciteA{Zheng2023} aplican esta metodología a planificación y programación de redes de
carga aérea. Su principal ventaja corresponde a la flexibilidad para explorar grandes
vecindarios, mientras que su principal limitación es la ausencia de una certificación de
optimalidad global.


% ============================================================
\subsubsection{Métricas de evaluación computacional}
\label{app:metricas}
% ============================================================

La comparación entre metodologías debe realizarse bajo condiciones computacionales
equivalentes.

Para los métodos exactos se consideran, como mínimo,

\[
\text{tiempo hasta la primera solución factible},
\]

\[
\text{mejor incumbente},
\qquad
\text{mejor bound},
\qquad
\text{brecha de optimalidad},
\]

\[
\text{tiempo total de resolución},
\qquad
\text{memoria utilizada}.
\]

Para un problema de maximización puede utilizarse

\begin{equation}
\mathrm{Gap}(\%)
=
100
\frac{UB-LB}
{\max\{1,|LB|\}},
\label{eq:gap}
\end{equation}

donde $LB$ representa el valor de la mejor solución factible conocida y $UB$ una cota
superior válida.

En generación de columnas se registrarán además el número de iteraciones, número de columnas
generadas y tiempo dedicado a los problemas de \textit{pricing}. Para Benders se considerarán
el número de iteraciones, número de cortes generados y evolución de las cotas.

En los métodos heurísticos y matheurísticos se comparará principalmente la calidad de la mejor
solución obtenida bajo un mismo presupuesto de tiempo.

Estas métricas permitirán identificar si la principal fuente de complejidad corresponde a la
secuenciación y programación temporal de las aeronaves, al flujo multicommodity de carga, al
número potencial de rotaciones o a la incorporación de restricciones operacionales
adicionales.